\chapter{Gute Form og dei ti boda}

\begin{motto}
Ti bod er nok til ein religion.\\ Dei er ikkje nok til ein vitskap.
\end{motto}

Den skandinaviske tradisjonen gav faget eit rykte for å vere godt. Den tyske etterkrigstradisjonen gav noko som ser ut til å vere meir: ein standard. Her, og berre her i heile historia, finn vi noko som liknar eit delt sett kriterium for kva godt design er, formulert, nummerert, skrive ned, og så breitt godteke at studentar frå Oslo til Seoul kan ramse det opp utanåt. Det er Dieter Rams sine ti prinsipp for godt design. Dei heng på veggen i atelier, dei vert siterte i jobbsøknader og festtalar, og dei er, utan samanlikning, det næraste designfaget nokon gong har kome ein felles fagleg målestokk. Difor er dette kapitlet det vanskelegaste i boka. Her må eg gå mot det beste motargumentet mot tesen min. Om faget har ein delt standard, manglar det vel ikkje eit fundament? Spørsmålet fortener eit ærleg svar, og svaret tek heile kapitlet, for det ligg gøymt i éin liten ting: dei ti boda kan ikkje ta feil.

Fyrst skal den tyske tradisjonen få det han skal ha, for han er den alvorlegaste i boka. Etter krigen stod Vest-Tyskland med eit øydelagt land og eit øydelagt namn, og industridesignen vart ein del av attreisinga, ikkje berre økonomisk, men moralsk. Paul Betts har vist korleis kvardagstinga bar ei tung byrde: dei skulle prove at landet kunne lage noko reint og ærleg att etter at alt anna var skitna til\autocite{betts2004}. «Gute Form», god form, var namnet, lansert mellom anna gjennom ei utstilling Max Bill sette saman i 1949, og det var meir enn ein stil. Det var ein moralsk kategori og eit oppdrag: tinga skulle vere sakssvarande, varige, fri for juks og staffasje, det motsette av nazismen sin patos og det motsette av den amerikanske stylingen sitt blendverk.

\section{Ulm og Braun}

Institusjonen er alt skildra: Hochschule für Gestaltung i Ulm, opna i 1953, arvtakaren til Bauhaus og den mest teorihungrige greina faget har sett, det miljøet som henta inn ergonomi, semiotikk og metodelære for å grunne form vitskapleg, og som leita etter fundamentet med all den strengskap tida åtte (kap. 8). Det vi treng her, er berre samarbeidet skulen gjekk inn i: Braun, der Hans Gugelot og den unge Dieter Rams, som kom dit i 1955, forma ein heil produktverd som vart sjølve biletet på Gute Form.

Sjå kor langt frå kvarandre ambisjonen og frukta enda. Det mest teoretisk ærgjerrige miljøet i fagets historie, det som henta inn informasjonsteori og semiotikk for å grunne form vitskapleg, leverte som sitt mest varige bidrag ei liste med ti setningar. Det er ikkje eit samantreff. Det er kva som hender når eit miljø vil ha ein teori, arbeider dugande i mange år, men ikkje har eit fundament å byggje teorien på. Det dei sit att med, er ei oppsummering av sin eigen gode smak, og dess klårare dei skriv henne ned, dess meir ser ho ut som ei lov.

\section{Objekta før orda}

For å sjå kva boda er, må ein sjå tinga dei skal oppsummere, og avstanden i tid mellom dei to. Rams kom til Braun i 1955, tjuetre år gamal, utdanna arkitekt og snikkar. Saman med Gugelot teikna han året etter platespelaren og radioen SK4: ein låg, kvit boks i metall og lakkert tre, reinsa ned til det minste, og over heile mekanismen eit lok av klårt pleksiglas, slik at platetallerken og tonearm og knappar låg openlyst til skode. Samtida kalla det \textit{Schneewittchensarg}, Snøkvits kiste, den gjennomsiktige kista der den vakre låg synleg under glaset. Folk greip etter eit eventyr for å forklare kjensla, og kjensla var rett: nokon hadde teke det som før vart pakka inn, og lagt det fram under glas. Loket var ikkje berre praktisk. Det var ein påstand. Denne maskinen lyg ikkje om kva han er, han gøymer ikkje mekanikken bak ein møbelfasade, han viser deg seg sjølv.

Men merk kva slags påstand det ikkje er. Det finst ingen grunn skriven ned for kvifor det openlyse er betre enn det innpakka, anna enn at det kjennest reinare og sannare. Ein annan smak, den borgarlege, som ville ha radioen som ein kommode mellom kommodane i stova, hadde like god rett til sitt, og SK4 har ingen teori som kan vise at han tek feil. Apparatet er framifrå formgjeve av to menn som meinte det. Det er berre ikkje eit prov på noko. T3-lommeradioen frå 1958 og lysbiletframvisaren D6 frå 1962 høyrer same verda til: ferdige, selde og elska, alle saman, lenge før det fanst ti bod.

Og det er den datoen som avgjer alt. Rams formulerte ikkje prinsippa på femtitalet då han byrja, og ikkje på sekstitalet då meisterverka kom. Han formulerte dei fyrst kring 1976, då han som vaksen designsjef vart beden om å seie kva som hadde styrt arbeidet hans gjennom tjue år. Sophie Lovell er utvitydig: prinsippa er ein retrospektiv destillasjon av ein praksis som alt var fullført, ikkje eit aksiomsett nokon praksis vart utleidd frå\autocite{lovell2011rams}. T3 vart ikkje laga etter regelen om at godt design er så lite design som mogleg; regelen vart skriven, tjue år seinare, av mannen som såg attende på T3. Boda er portrettet, ikkje oppskrifta, og Snøkvits kiste er andletet portrettet vart måla etter.

\section{Dei ti boda lese som det dei er}

Lat oss då lese dei. Godt design er nyskapande. Godt design gjer eit produkt nyttig. Godt design er estetisk. Godt design gjer eit produkt forståeleg. Godt design er diskret. Godt design er ærleg. Godt design er varig. Godt design er konsekvent ned til siste detalj. Godt design er miljøvenleg. Og det tiande, det mest siterte av alle: godt design er så lite design som mogleg, \textit{Weniger, aber besser}, mindre, men betre.

Les dei ein gong til og spør kva slags utsegner dette er. Definisjonar er det ikkje; dei føreset alle at vi alt veit kva godt design er, og listar opp eigenskapar det har. Målbare kriterium er det ikkje; kor mykje er «så lite som mogleg», når er noko diskret nok, kvar går grensa mellom det nødvendige og det overflødige? Men det avgjerande spørsmålet er det tredje: kan dei vere usanne? Tenk deg eit godt, varmt, elska objekt som er rikt dekorert, mangefarga og overdådig, det motsette av diskret. Har Rams då teke feil? Nei. Tilhengaren seier berre at objektet ikkje er godt design, fordi godt design er pyntelaust. Bodet kan ikkje motseiast av noko døme, av di det definerer bort kvart moteksempel på førehand.

Det hjelper å sjå kva ein påstand som faktisk kan vere usann, er for noko. Når ein ingeniør seier at ein bjelke ber, kan verda svare nei: bjelken brest under lasta. Påstanden har eit innhald som rekk ut over den som ytrar han, og difor finst det eit utfall som ville motsagt han. Aforismen har ikkje slikt innhald. Når Rams seier at ei form er god fordi ho er pyntelaus, finst det inga last ho kan brest under, ikkje av di design er ein mislukka ingeniørdisiplin, men av di smak heva til lov er eit anna slag ting enn ein prøvbar påstand. Kvart pyntelaust ting som mislukkast, vert forklart som dårleg utført pyntelausheit; kvart rikt ting som lukkast, som berre populært, ikkje godt. Prinsippet er immunt, og «mindre, men betre» grunngjev ikkje eingong seg sjølv: det føreset det det skulle prove, at mindre er betre. Det er sirkelen all smak går i når han vil bli til lov utan å gå vegen om ein teori. Eit rikt, elska objekt står utanfor sirkelen, ikkje motsagt, berre definert bort.

\section{Frå maksime til marknadsføring}

Eit prinsipp som ikkje kan grunngje, måle eller motseie seg sjølv, har ein farleg eigenskap: det kan tileignast av kven som helst som likar resultatet, gratis. Og det vart det. Den mest kjende arvtakaren etter Rams er ikkje ein teori, men eit selskap. Jony Ive, som leidde formgjevinga ved Apple frå nittitalet, har sjølv ope vedkjent gjelda til Braun, og ho er sporbar ned til knappen. Set Braun-rekneskapsmaskina ET66 frå 1977, teikna av Rams og Dietrich Lubs, ved sida av kalkulatorappen på den fyrste iPhone frå 2007: same runde knappar med svak kvelving, same fargekode med likskapsteiknet i oransje, same haldne klårleik. Set T3 ved sida av iPod, som Ive sjølv har kalla ein etterkomar. Dette er ikkje vag inspirasjon. Det er eit nært studium av ein konkret formverd, ført over frå tyske kvitevarer til amerikansk forbrukselektronikk.

Men sjå kva som kryssa og kva som ikkje gjorde det. Forma kryssa: dei reine flatene, fråvêret av dekor, paletten, «mindre, men betre» som visuelt program. Grunngjevinga kryssa ikkje. Hjå Rams var den nakne forma bunden til at tinget skulle vere ærleg, sakssvarande og varig; bod sju seier rett ut at godt design er langlivd, bod ni at det er miljøvenleg. Apple tok forma og let dei to boda liggje att på golvet. Produkta vart bundne til lukka system, lagde for å skiftast ut i takt med årlege lanseringar, med batteri som var vanskelege å byte og innkapsla bak limte deksel, sjølve den planlagde foreldinga Rams-etikken var meint å stå imot. Forma sa framleis «ærleg og varig» med Rams si stemme. Produktet praktiserte det motsette, og målestokken merka ingen skilnad.

Hadde Rams gjeve oss eit «fordi», ein grunn til at mindre er betre, så ville den grunnen reist seg mot lånet: nei, ein telefon som ikkje kan opnast og er meint å døy om tre år, er ikkje godt design etter mine eigne premiss, for varigheit var grunnen, ikkje pynten. Men aforismen har ikkje noko «fordi». Han har eit utsjånad, og eit utsjånad kan kopierast utan grunnen sin og festast på det motsette av det grunnen ville tillate. Ein kan innvende at Apple sjølv ville kalle reinleiken ein etikk, ein respekt for brukaren. Men nett den innvendinga viser kor lite boda ber, for reinleik kan like gjerne grunngjevast som luksussignal, som distinksjon, som ein måte å gjere det dyre til det moralsk overlegne, og boda har ikkje noko språk til å skilje dei to bruksmåtane. Donald Norman har minna oss om at den emosjonelle bindinga vår til ting aldri vart fanga av reinleiken åleine\autocite{norman2004emotional}; her er poenget skarpare endå. Den runde, klåre knappen frå ET66 vandra fritt frå ein maskin laga for å reknast med i tjue år, inn på ein skjerm som etterliknar han, og bar etikkens utsjånad utan etikkens innhald.

\section{Laget bak mannen}

Aforismen om Rams løyner ein ting som fortel mykje om kva slags arbeid Gute Form eigenleg var. Han er individuell, \textit{Rams} sine ti bod, \textit{Rams} sin smak, \textit{Rams} si hand, og han matar ein heltekultus der éin mann ser sanninga om form og skriv henne ned. Men Braun-tinga vart laga av eit lag. Gerd Alfred Müller, på Braun frå midt-femtitalet, teikna seinare, etter at han gjekk, den sekskanta Lamy-pennen som framleis vert seld; det beviset på at smaken låg i handa og ikkje i Rams sine bod, han teikna ho aldri. Reinhold Weiss teikna vifteomnen H1, Fritz Eichler styrte den kulturelle retninga, Lubs teikna klokker og talloppsett. Rams var sjefen og det fremste hovudet, men han var ikkje laget.

Det laget makta, var ikkje ein teori. Det var ein husstil. Mange ulike hender produserte produkt som høyrde saman, som såg ut som søsken, slik at ein barbermaskin og ein radio og ein kjøkenmaskin frå same hus straks kjentest att som same familie. Dette er den same tause kunnskapen vi byrja boka med (kap. 1): ein \textit{knowing how} og ikkje ein \textit{knowing that}, ein innøvd, delt smak ført frå hovud til hand i eit konkret miljø, som ikkje bur i ord men i augo og hendene til folk som har arbeidd lenge nok saman til å sjå likt. Han kan ikkje skrivast ned uttømmande, fordi han ikkje er ein setning. Braun-familien er det beste argumentet i heile boka for at god form finst og kan delast, og samstundes det beste mot at ho er ein teori, for det ho vart delt som, var ein praksis ein måtte stå inni for å eige. Og her ligg ironien: faget tok det individuelle ut av Braun, mannen, boda, plakaten, og let det kollektive liggje, fordi det individuelle let seg sitere og det kollektive ikkje gjer det. Eit lag sin husstil kan du ikkje hengje på veggen. Du kan berre lære han ved å vere der.

\section{«Eg ville ikkje blitt designar i dag»}

Det finst eit vitne mot dei ti boda som er vanskelegare å avvise enn nokon kritikar, og det er Rams sjølv. Mannen som gav faget sitt næraste til ein felles standard, vart med åra ein motvillig profet. Som godt vaksen uttrykte han uro over kor mykje verda no produserer, over flaumen av ting, over forbruket boda skulle vore eit vern mot og som i staden berre voks. Og han sa, i ei utsegn som har gått verda rundt og som vert ståande att i Gary Hustwits dokumentar frå 2018, at om han skulle velje på nytt, ville han ikkje blitt designar i dag.

Kva fortvilar han over? Ikkje at han tok feil om kva god form er; han stod ved smaken sin til det siste. Han fortvilar over at det ikkje hjelpte. Bod sju sa at god form er varig, bod ni at ho er miljøvenleg, bod ti at ho er så lite som råd, og likevel levde Rams lenge nok til å sjå ein verd meir overfylt av kortlivde, lygande ting enn nokon gong, og til å sjå sitt eige reine formspråk pryde mange av desse tinga som eit kvalitetsstempel. Han hadde gjeve faget eit motgift, og motgiftet var blitt ein etikett.

Men sjå kva denne desillusjonen var, for forma hans stadfestar diagnosen. Rams konkluderte ikkje med at boda var feilkonstruerte, at dei mangla ein teori, at dei var ein smak heva til lov. Han konkluderte moralsk og personleg, med sorg. Når ein teori sviktar, kan ein peike på kvar han svikta og rette han; det er ein fagleg operasjon som fører kunnskapen vidare. Når ein smak ser at verda ikkje vart som han håpa, har han ingen stad å rette noko. Rams hadde ikkje noko språk til å seie «boda var gale på dette punktet, difor kunne dei misbrukast slik», fordi boda aldri var den slags ting ein kan ha rett eller gale om. Den som hadde ein teori, ville feilsøkt. Rams, som hadde ein smak, kunne berre sørgje og trekkje seg attende.

\section{Den udøyelege aforismen}

Rams sine ti prinsipp er det beste faget har. Dei er klokare, meir samanhengande og meir ærlege enn nesten alt anna i denne boka, og ein student som tok dei på alvor ville lage betre ting enn ein som ikkje gjorde det. Innvendinga er ikkje at dei er dårlege. Innvendinga er at vi har forveksla dei med eit fundament. Vi har late som om ti aforismar frå ein framifrå formgjevar er det same som ein teori om form, at det å kunne ramse dei opp er å vite kvifor ei form verkar, og i den forvekslinga har vi gjort fundamentet faget treng vanskelegare å sjå, ikkje lettare. For når nokon spør om designfaget har ein delt standard, peikar vi på boda og er nøgde. Men ein standard du ikkje kan måle med, ikkje grunngje og ikkje motseie, er ikkje ein standard. Det er ei truvedkjenning. Og ei truvedkjenning kan samle ei rørsle. Ho kan ikkje bere ein vitskap.

Det skal seiast at boda gjer ein ærleg ting dei vakre orda i Norden ikkje gjorde: dei prøver i det minste å vere eksplisitte. Dei skriv ned kriteria i staden for å la dei liggje i smaken til den som dømer. Det er eit framsteg, og det peikar mot noko. For om du fyrst har byrja å skrive ned kva som gjer form god, er du eit kort steg frå det neste og farlegare spørsmålet: kan dette gjerast til ein metode? Kan ein skrive ikkje berre kva godt design er, men korleis ein kjem fram til det, steg for steg, slik at det ikkje lenger heng på handa og auget til meisteren? Nokon stilte nett det spørsmålet, og dei stilte det med eit alvor og ein optimisme faget aldri har sett maken til. Dei kalla det design methods.

\laeringsmal{\begin{itemize}
\item Gjere greie for «Gute Form» og Braun--Ulm-koplinga, og for Rams' ti prinsipp i historisk kontekst.
\item Skilje ein eksplisitt aforisme frå ein prøvbar påstand, og forklare kvifor det fyrste ikkje er det andre.
\item Drøfte korleis dei ti boda seinare vart marknadsføring (Apple/Ive-arven).
\end{itemize}}

\ovingar{\begin{enumerate}
\item Vel eitt av Rams' ti bod og prøv å operasjonalisere det: korleis ville du \emph{måle} om ei gjeven form oppfyller det? Kvar bryt forsøket saman?
\item «Godt design er så lite design som mogleg.» Er dette ein målestokk eller ein smak? Grunngje.
\item Set boda ved sida av eit kriteriesett frå eit fag som har eit etisk organ, til dømes ei klinisk retningslinje. Poenget er ikkje at design burde vore medisin, men at samanlikninga gjer synleg kva slags utsegn boda er, og kva dei ikkje er. Kva kan retningslinja gjere som boda ikkje kan?
\end{enumerate}}

\vidarelesing{\fullcite{lovell2011rams}; \fullcite{betts2004}; \fullcite{lindinger1991ulm}; \fullcite{norman2004emotional}}
